\documentclass[10pt,a4paper]{amsart}
\usepackage{amsmath,amssymb,amsthm,graphicx}
% Retained minimal preamble for the Stacks Project chapter more-algebra.tex (More on Algebra).
% The chapter file pinned-original-source.tex carries only \input{preamble} on its first line, so no
% preamble text lives inside the chapter itself.  The Stacks Project uses one single shared preamble for
% every chapter; the snapshot stages it identically as sources/stacks-chow/preamble.tex and
% sources/stacks-etale-cohomology/preamble.tex.  The block below is copied mechanically, line for line,
% from sources/stacks-chow/preamble.tex lines 14-18 (the Xy-pic setup) and 154-217 (the native theorem
% declarations together with the mathematical macros).  Nothing was added, removed or redefined; the
% wrapper class and the amsmath/amssymb/amsthm/graphicx base packages are supplied by the assembler.
% Chapter: more-algebra.tex -- More on Algebra.
% For commutative diagrams we use Xy-pic
\usepackage[all]{xy}

% We use 2cell for 2-commutative diagrams.
\xyoption{2cell}

% Theorem environments.
%
\theoremstyle{plain}
\newtheorem{theorem}[subsection]{Theorem}
\newtheorem{proposition}[subsection]{Proposition}
\newtheorem{lemma}[subsection]{Lemma}

\theoremstyle{definition}
\newtheorem{definition}[subsection]{Definition}
\newtheorem{example}[subsection]{Example}
\newtheorem{exercise}[subsection]{Exercise}
\newtheorem{situation}[subsection]{Situation}

\theoremstyle{remark}
\newtheorem{remark}[subsection]{Remark}
\newtheorem{remarks}[subsection]{Remarks}

\numberwithin{equation}{subsection}

% Macros
%
\def\lim{\mathop{\mathrm{lim}}\nolimits}
\def\colim{\mathop{\mathrm{colim}}\nolimits}
\def\Spec{\mathop{\mathrm{Spec}}}
\def\Hom{\mathop{\mathrm{Hom}}\nolimits}
\def\Ext{\mathop{\mathrm{Ext}}\nolimits}
\def\SheafHom{\mathop{\mathcal{H}\!\mathit{om}}\nolimits}
\def\SheafExt{\mathop{\mathcal{E}\!\mathit{xt}}\nolimits}
\def\Sch{\mathit{Sch}}
\def\Mor{\mathop{\mathrm{Mor}}\nolimits}
\def\Ob{\mathop{\mathrm{Ob}}\nolimits}
\def\Sh{\mathop{\mathit{Sh}}\nolimits}
\def\NL{\mathop{N\!L}\nolimits}
\def\CH{\mathop{\mathrm{CH}}\nolimits}
\def\proetale{{pro\text{-}\acute{e}tale}}
\def\etale{{\acute{e}tale}}
\def\QCoh{\mathit{QCoh}}
\def\Ker{\mathop{\mathrm{Ker}}}
\def\Im{\mathop{\mathrm{Im}}}
\def\Coker{\mathop{\mathrm{Coker}}}
\def\Coim{\mathop{\mathrm{Coim}}}

% Boxtimes
%
\DeclareMathSymbol{\boxtimes}{\mathbin}{AMSa}{"02}

%
% Macros for moduli stacks/spaces
%
\def\QCohstack{\mathcal{QC}\!\mathit{oh}}
\def\Cohstack{\mathcal{C}\!\mathit{oh}}
\def\Spacesstack{\mathcal{S}\!\mathit{paces}}
\def\Quotfunctor{\mathrm{Quot}}
\def\Hilbfunctor{\mathrm{Hilb}}
\def\Curvesstack{\mathcal{C}\!\mathit{urves}}
\def\Polarizedstack{\mathcal{P}\!\mathit{olarized}}
\def\Complexesstack{\mathcal{C}\!\mathit{omplexes}}
% \Pic is the operator that assigns to X its picard group, usage \Pic(X)
% \Picardstack_{X/B} denotes the Picard stack of X over B
% \Picardfunctor_{X/B} denotes the Picard functor of X over B
\def\Pic{\mathop{\mathrm{Pic}}\nolimits}
\def\Picardstack{\mathcal{P}\!\mathit{ic}}
\def\Picardfunctor{\mathrm{Pic}}
\def\Deformationcategory{\mathcal{D}\!\mathit{ef}}

% --- appended 2026-10-12 (Chief): the project own definitions of the reference/slogan/history
% display environments, copied mechanically from the shared Stacks preamble lines 8-12
% (\usepackage{verbatim} plus the three \newenvironment lines, which the chapter preamble keeps
% verbatim from the project preamble). Nothing added, removed or redefined: these environments
% carry no mathematical content and the project itself typesets them as comments. Without these
% three lines a counted statement whose own text uses them fails with Environment slogan undefined.
\usepackage{verbatim}
\newenvironment{reference}{\comment}{\endcomment}
\newenvironment{slogan}{\comment}{\endcomment}
\newenvironment{history}{\comment}{\endcomment}
\usepackage{cleveref}
\crefname{originallemma}{original lemma}{original lemmas}
\Crefname{originallemma}{Original lemma}{Original lemmas}
\makeatletter
\newlabel{algebra-lemma-Noetherian-power-ideal-kills-module}{{00L6}{source}{}{sourcepointer-1}{}}
\newlabel{algebra-lemma-Noetherian-power-ideal-kills-module@cref}{{[originallemma][0][]00L6}{[1][0][]source}{}{}{}}
\makeatother
\crefname{originallemma}{original lemma}{original lemmas}
\Crefname{originallemma}{Original lemma}{Original lemmas}
\makeatletter
\newlabel{algebra-lemma-snake}{{07JW}{source}{}{sourcepointer-2}{}}
\newlabel{algebra-lemma-snake@cref}{{[originallemma][0][]07JW}{[1][0][]source}{}{}{}}
\makeatother
\crefname{originallemma}{original lemma}{original lemmas}
\Crefname{originallemma}{Original lemma}{Original lemmas}
\makeatletter
\newlabel{algebra-lemma-support-point}{{00L5}{source}{}{sourcepointer-3}{}}
\newlabel{algebra-lemma-support-point@cref}{{[originallemma][0][]00L5}{[1][0][]source}{}{}{}}
\makeatother
\crefname{originallemma}{original lemma}{original lemmas}
\Crefname{originallemma}{Original lemma}{Original lemmas}
\makeatletter
\newlabel{lemma-polypoly}{{37979}{source}{}{sourcepointer-4}{}}
\newlabel{lemma-polypoly@cref}{{[originallemma][0][]37979}{[1][0][]source}{}{}{}}
\makeatother
\title{In a one-dimensional Noetherian local ring the number of minimal primes is bounded by the length of the quotient by an element avoiding them}
\author{The Stacks Project Authors}
\date{stacks-project:more-algebra.tex:local-snapshot-20261012; GNU FDL 1.2 or later}
\begin{document}
\maketitle
\setcounter{subsection}{1}
\noindent\textit{Editorial note.} All delivered mathematical text is the author's own, copied from the pinned Stacks Project chapter file more-algebra.tex; no statement or proof step was rewritten, reordered or removed.\par\medskip
\noindent\textit{Editorial note.} The chapter file pinned-original-source.tex carries only \textbackslash{}input\{preamble\} on its first line, so no preamble text lives inside the chapter. The retained file sources/stacks-more-algebra/retained-preamble.tex is a mechanical, line-for-line copy of the project's single shared preamble as staged in the snapshot; nothing was added, removed or redefined and no package was added.\par\medskip
\noindent\textit{Editorial note.} The source URL is the source-level Stacks Project URL, as required when more than one card is produced from one source; the per-claim tag is carried by claim\_key.\par\medskip
\noindent\textit{Editorial note.} The author writes 'as there is a surjection R/xR -> R'/xR'' and concludes the length of the target is smaller; the conclusion as stated is the intended one and is reproduced unchanged, as is the phrase 'with minimal primes p\_1, ..., p\_t' following 'Assume R is reduced'.\par\medskip
\noindent\textit{Allowed standard prerequisites (editorial).}\begin{itemize}
\item Support of a finite module and finite length for a module supported at the closed point. used at proof lines 38048-38052. Shows Q has finite length and hence is annihilated by a power of x.
\item The snake lemma applied to a square of multiplication maps. invoked at proof lines 38067-38070. Produces the exact sequence relating R/x\textasciicircum{}nR, M/x\textasciicircum{}nM and Q, which yields the length equality.
\item The one-dimensional length lemma for the quotients R\_i/xR\_i. invoked at proof lines 38076-38077. Converts the length equality into the explicit alternating count from which the bound follows.
\end{itemize}\medskip
\ifdefined\hypertarget\hypertarget{sourcepointer-1}{}\fi
\noindent\textit{Source reference.} Original lemma 00L6: focused-seeds/AGP/sources/stacks-snapshot-20261011/tags.txt, source lines 776--776. Cross-chapter navigation pointer into the Stacks Project; the standard fact it supplies is applied inside the counted proof span and the argument resting on it stays there.\par\medskip
\ifdefined\hypertarget\hypertarget{sourcepointer-2}{}\fi
\noindent\textit{Source reference.} Original lemma 07JW: focused-seeds/AGP/sources/stacks-snapshot-20261011/tags.txt, source lines 9636--9636. Cross-chapter navigation pointer into the Stacks Project; the standard fact it supplies is applied inside the counted proof span and the argument resting on it stays there.\par\medskip
\ifdefined\hypertarget\hypertarget{sourcepointer-3}{}\fi
\noindent\textit{Source reference.} Original lemma 00L5: focused-seeds/AGP/sources/stacks-snapshot-20261011/tags.txt, source lines 775--775. Cross-chapter navigation pointer into the Stacks Project; the standard fact it supplies is applied inside the counted proof span and the argument resting on it stays there.\par\medskip
\ifdefined\hypertarget\hypertarget{sourcepointer-4}{}\fi
\noindent\textit{Source reference.} Original lemma 37979: sources/stacks-more-algebra/pinned-original-source.tex, source lines 37979--37979. Same-chapter navigation pointer; the statement it names is standard for the section and the argument that consumes it stays inside the counted proof span.\par\medskip
\begin{lemma}
\label{lemma-minprimespoly}
Let $(R, \mathfrak m)$ be a Noetherian local ring of dimension $1$.
Let $x \in \mathfrak m$ be an element not contained in any minimal
prime of $R$. Let $t$ be the number of minimal prime ideals of $R$.
Then $t \leq \text{length}_R(R/xR)$.
\end{lemma}

\begin{proof}
Let $\mathfrak p_1, \ldots, \mathfrak p_t$ be the minimal prime ideals
of $R$. Set $R' = R/\sqrt{0} = R/(\bigcap_{i = 1}^t \mathfrak p_i)$.
We claim it suffices to prove the lemma for $R'$. Namely, it is clear
that $R'$ has $t$ minimal primes
too and $\text{length}_{R'}(R'/xR') = \text{length}_R(R'/xR')$
is less than $\text{length}_R(R/xR)$ as there is a surjection
$R/xR \to R'/xR'$. Thus we may assume $R$ is reduced.

\medskip\noindent
Assume $R$ is reduced with minimal primes
$\mathfrak p_1, \ldots, \mathfrak p_t$.
This means there is an exact sequence
$$
0 \to R \to
\prod\nolimits_{i = 1}^t R/\mathfrak p_i \to Q \to 0
$$
Here $Q$ is the cokernel of the first map.
Write $M = \prod_{i = 1}^t R/\mathfrak p_i$.
Localizing at $\mathfrak p_j$ we see that
$$
R_{\mathfrak p_j} \to M_{\mathfrak p_j} =
\left(\prod\nolimits_{i=1}^t R/\mathfrak p_i\right)_{\mathfrak p_j} =
(R/\mathfrak p_j)_{\mathfrak p_j}
$$
is surjective. Thus $Q_{\mathfrak p_j} = 0$ for all $j$.
We conclude that $\text{Supp}(Q) = \{\mathfrak m\}$ as $\mathfrak m$
is the only prime of $R$ different from the $\mathfrak p_i$.
It follows that $Q$ has finite length
(Algebra, Lemma \ref{algebra-lemma-support-point}).
Since $\text{Supp}(Q) = \{\mathfrak m\}$ we
can pick an $n \gg 0$ such that $x^n$
acts as $0$ on $Q$
(Algebra, Lemma \ref{algebra-lemma-Noetherian-power-ideal-kills-module}).
Now consider the diagram
$$
\xymatrix{
0 \ar[r] & R \ar[r] \ar[d]^-{x^n} & M
\ar[r] \ar[d]^-{x^n} & Q \ar[r] \ar[d]^-{x^n} & 0 \\
0 \ar[r] & R \ar[r] & M \ar[r] & Q \ar[r] & 0
}
$$
where the vertical maps are multiplication by $x^n$. This is injective on
$R$ and on $M$ since $x$ is not contained in any of the $\mathfrak p_i$.
By the snake lemma (Algebra, Lemma \ref{algebra-lemma-snake}),
the following sequence is exact:
$$
0 \to Q \to R/x^nR \to M/x^nM \to Q \to 0
$$
Hence we find that $\text{length}_R(R/x^nR) = \text{length}_R(M/x^nM)$
for large enough $n$. Writing $R_i = R/\mathfrak p_i$ we see
that $\text{length}(M/x^nM) =
\sum_{i = 1}^t \text{length}_R(R_i/x^nR_i)$.
Applying Lemma \ref{lemma-polypoly} and the fact that $x$ is a nonzerodivisor
on $R$ and $R_i$, we conclude that
$$
n \text{length}_R(R/xR) =
\sum\nolimits_{i = 1}^t n \text{length}_{R_i}(R_i/x R_i)
$$
Since $\text{length}_{R_i}(R_i/x R_i) \geq 1$ the lemma is proved.
\end{proof}

\end{document}
